% Part of the present nc cubic threefold moduli manuscript.
% Restored from the 2026-11-01 central-pair paper V, sections 19--21.
\section{全混合族の低次数中心、偶数Veroneseと剛体ambient}
\label{sec:moduli-rigid-ambient}
本節では、可換点の全Artin比較とは別に、有限非可換パラメータ
において一つのambientを固定したまま三次方程式を動かす問題を調べる。
元の中心三次対原稿にある命題を現在有効な仮定と証明に整理する。
非可換Bertiniの適用、三次Sklyaninの位数条件、
および有限整数証明書の独立監査は未完了である。

\subsection{三次商による二次中心の非増大}
$U=\Span(u,v)$、$R=k[u,v]$、$k=\C$とおく。
$F\in k[x,y,z,u,v]_3$は
$F(x,y,z,0,0)$が滑らかな平面三次式$f_3(x,y,z)$であると仮定する。
$h\in k^\times$に対して
\[
 B(h,F)=\frac{R\langle x,y,z\rangle}
 {([x,y]-h\Symm F_z,\ [y,z]-h\Symm F_x,\
    [z,x]-h\Symm F_y)}
\]
とし、次数三の中心正則元$W$で割った
$A=B/(W)$を考える。ここで$B$は一次関係を持たないとする。

\begin{lemma}[中心三次商では二次中心が増えない]
\label{lem:moduli-no-new-quadratic-center}
上の仮定で、自然な同一視$B_2\simeq A_2$の下で
$Z(A)_2=Z(B)_2$となる。
\end{lemma}
\begin{proof}
$Q\in B_2$の像が$A$で中心なら、
$[Q,x_i]=c_iW$、$c_i\in k$、
$(x_1,x_2,x_3)=(x,y,z)$と書ける。
$u,v,W$が中心なので
\[
 0=[Q,[x_i,x_j]]
 =h[Q,\Symm F_{x_k}]
 =hW\sum_\ell c_\ell F_{x_kx_\ell}
\]
である（$i,j,k$は巡回）。
$hW$を消去すると、右側の一次式が$B_1$で零。
一次関係を持たないから可換多項式として零であり、
$u=v=0$で
\[
  \sum_\ell c_\ell(f_3)_{x_kx_\ell}=0
  \quad(k=1,2,3).
\]
すなわち$\sum_\ell c_\ell(f_3)_{x_\ell}$は
定数の斉次二次式なので零になる。
$c\ne0$なら座標を取り換えれば$f_3$は一変数を含まず、
対応する射影点で偏微分がすべて消えて滑らかさに矛盾する。
従って$c=0$で、$Q$は既に$B$の中心である。
\end{proof}

\begin{lemma}[全混合族の次数一・二の中心]
\label{lem:moduli-low-center}
$S=B/(u,v)$とし、中心三次元$\gamma\in S_3$で
\[
 S/(\gamma)\simeq B(E,\mathcal L,\tau)
\]
となるとする。$E$は滑らかな楕円曲線、
$\deg\mathcal L=3$、$\tau$は平行移動で$\tau^2\ne1$と仮定する。
さらに相対PBWにより
$\gr_{x,y,z} B\simeq S\otimes_k k[u,v]$とする。
このとき
\[
 Z(B)_1=Z(A)_1=U,\qquad
 Z(B)_2=Z(A)_2=\Sym^2U.
\]
\end{lemma}
\begin{proof}
twisted homogeneous coordinate ringは
$k(E)[t,t^{-1};\tau]$に埋め込まれる。
そこに次数$m=1,2$の中心元$at^m$があれば
一次切断$s_1,s_2$との可換性から
$\tau^m(s_1/s_2)=s_1/s_2$を得る。
次数三の完全線形系の一次切断の比は$k(E)$を生成するため
$\tau^m=1$となり仮定に反する。
従って$S/(\gamma)$には低次数中心がなく、
$S\to S/(\gamma)$は次数一・二で単射だから
$Z(S)_1=Z(S)_2=0$。
$B$の中心元の最高相対次数を取ると、
正の相対次数一・二は消え、中心は$k[u,v]$に属する。
三次商でも次数一の関係は増えず、
二次は補題\ref{lem:moduli-no-new-quadratic-center}を使う。
\end{proof}

\subsection{偶数Veroneseからの標識付き代数の回復}
\begin{lemma}[中心二変数の平方根と奇数部分]
\label{lem:moduli-even-veronese}
$A$はconnectedな次数付き$k$代数で、
二つの一次中心元$u,v$が多項式環$S=k[u,v]$を生成し、
$A$が$S$上自由、
$Z(A)_2=S_2$とする。
$R=A^{(2)}=\bigoplus_{n\ge0}A_{2n}$とおくと、
$R$の次数付き同型類は$A$の次数付き同型類を決定する。
詳しくは、同じ仮定を満たす$A'$について、
任意の$\phi:R\xrightarrow{\sim}(A')^{(2)}$は
正確に二つの次数付き同型$A\to A'$へ持ち上がり、
その差は奇数次数に$-1$を掛ける自己同型である。
\end{lemma}
\begin{proof}
$z\in Z(A)_1$について$uz,vz\in S_2$。
$u(vz)=v(uz)$と$S$の一意分解性から
$uz=u\ell$、$\ell\in S_1$であり、
$u$の正則性から$z=\ell$となる。
$z\in Z(R)$と奇数次数の$a\in A$に対しては
$ua\in R$なので
$u[z,a]=[z,ua]=0$より$[z,a]=0$。
従って$Z(R)=Z(A)^{(2)}$。
その次数一中心の生成する内在的部分環は
\[
 D=k[u^2,uv,v^2]
   \simeq k[\alpha,\beta,\gamma]/(\alpha\gamma-\beta^2).
\]
$\Proj D$は二次曲線で$\OO_{\Proj D}(1)=\OO_{\PP^1}(2)$。
任意の$D\simeq D'$は二次Veronese表現
$\GL_2\to\Aut_{\gr}D$により二通りだけ
$S\simeq S'$へ持ち上がり、その差は全一次元の符号反転。

$u,v$を局所化して$R$の次数を二倍すると
\[
 A[u^{-1}]
  \simeq R[(u^2)^{-1}]
    \otimes_{D[(u^2)^{-1}]}S[u^{-1}]
\]
であり、$v$の場合も同様である。
$A$は$S$上自由なので
\[
 A[u^{-1}]\cap A[v^{-1}]=A
 \quad\text{in }A[(uv)^{-1}].
\]
これが$R$と$S$だけからの復元式である。
同型を持ち上げた後、偶数次数と$S_1$上で恒等なら、
奇数元$a$について$u\phi(a)=\phi(ua)=ua$から
$\phi(a)=a$。残る一つは全奇数次数の符号反転である。
\end{proof}

\begin{theorem}[構造対象を保つ導来同値：条件付き全混合復元]
\label{thm:moduli-marked-derived}
$A=B/(W)$と$A'=B'/(W')$は上の全混合族の中心正則三次商で、
補題\ref{lem:moduli-low-center}の仮定と
それぞれの$k[u,v]$上の自由性を満たすとする。
両方の$\qgr$は有限大域次元を持つと仮定する。
次は同値である：
\begin{enumerate}
\item $A\simeq_{\gr}A'$;
\item $\OO_A$を$\OO_{A'}$に送る$k$線形Abel圏同値
 $\qgr A\simeq\qgr A'$が存在する;
\item ある$m,r\in\mathbb Z$について
 $\OO_A$を$\OO_{A'}(m)[r]$に送る$k$線形三角同値
 $D^b(\qgr A)\simeq D^b(\qgr A')$が存在する。
\end{enumerate}
これらの条件は、周囲の二次代数と三次中心直線
$(B,kW)$の同型類も決定する。
\end{theorem}
\begin{proof}
$(1)\Rightarrow(2)\Rightarrow(3)$は標準的である。
$(3)$の同値にshiftとtwistを合成して構造対象を一致させる。
$B$はCY五次元、$W$は中心正則三次元だから
balanced dualizing complexは$A(-2)[4]$である。
tailsの有限大域次元を使ってSerre関手は
$\mathsf S_A\simeq(-2)[3]$となる。
これは\href{https://www.math.bgu.ac.il/~amyekut/publications/nc-serre/paper.pdf}
{Yekutieli--ZhangのSerre双対}の適用を入力とする。
同値はSerre関手と可換なので
$\mathsf H_A=\mathsf S_A^{-1}[3]\simeq(2)$を保ち、
その軌道環
\[
 \bigoplus_{n\ge0}\Hom(\OO_A,\mathsf H_A^n\OO_A)
 \simeq A^{(2)}
\]
を積とともに同定する。
補題\ref{lem:moduli-low-center}と
\ref{lem:moduli-even-veronese}により$A\simeq A'$。
最後に二次関係の核から$B$を、三次の核から$kW$を回復する
定理\ref{thm:moduli-iso}を使う。
\end{proof}
\begin{remark}
\label{rem:moduli-marking}
ここで除いたのは次数一shiftの\emph{指定}であり、
構造対象の条件自体は残る。任意の導来自己同値が
構造対象をline objectのshiftへ送るとは主張しない。
有限大域次元を持たないtailsへの適用も主張しない。
\end{remark}

\subsection{全混合ポテンシャルの有限同型判定}
\begin{proposition}[中心blockと三次ポテンシャルの共変性]
\label{prop:moduli-mixed-block-isomorphism}
上の低次数中心を持つ$B(h,F)$、$B(h',F')$について
任意の次数付き同型$B(h',F')\to B(h,F)$は
\[
 s' \longmapsto Ns,\qquad x'\longmapsto Lx+Ms,
 \quad
 L\in\GL_3(k),\ N\in\GL_2(k),\ M\in\mathrm{Mat}_{3\times2}(k)
\]
の形を持つ。これは同型である必要十分条件が
ある$r\in k[u,v]_3$について
\[
 h'F'(Lx+Ms,Ns)=h\det(L)F(x,s)+r(s)
\]
となることに等しい。
中心商の次数付き同型にはさらに
$\psi_{L,M,N}(W')=\kappa W$ ($\kappa\in k^\times$)
を要し、これで必要十分である。
\end{proposition}
\begin{proof}
低次数中心が$U=\Span(u,v)$なので同型はblock形となる。
$([y,z],[z,x],[x,y])^t$の変換は
$\det(L)L^{-t}$であり、関係式の保存条件は
\[
 h'\nabla_{x'}F'(Lx+Ms,Ns)
 =h\det(L)L^{-t}\nabla_xF(x,s).
\]
二次関係空間が交代二次式から対称二次式へのgraphであるため、
これは可換二次多項式の恒等式と同値。
chain ruleで勾配の式を積分すると、差は$k[u,v]_3$に入る。
最後の中心直線保存は二次包絡と三次核の回復に等しい。
\end{proof}

\subsection{射影自己同型のない明示ambient}
次の例は元の中心三次対原稿第21節にあったものであり、
正則な中心商の開集合の存在は、中心商論文の有限位数境界に関する定理から従う。

\begin{proposition}[明示的な剛体ambient：中心の入力を指定]
\label{prop:moduli-rigid-example}
\[
 F_*=\frac{x^3+y^3+z^3}{3}
     +uxy+vyz+(u+v)zx
\]
に対する$h=1$の代数
\[
 B_*=\frac{k[u,v]\langle x,y,z\rangle}
 {([x,y]-z^2-vy-(u+v)x,\
   [y,z]-x^2-uy-(u+v)z,\
   [z,x]-y^2-ux-vz)}
\]
を考える。$S_*=B_* /(u,v)$の中心三次商が、
平面三次$E_*:x^3+y^3+z^3+xyz=0$に付随する
位数6の平行移動を持つTHCRであることを入力とする。
このとき
\[
 Z(B_*)_1=ku\oplus kv,\qquad
 Z(B_*)_3=kG_*\oplus k[u,v]_3,
 \qquad \Aut_{\gr}(B_*)=k^\times.
\]
ここで
\[
\begin{aligned}
 \Omega_1&=\frac{x^3+y^3+z^3}{3}
       -\frac19\sum_{\sigma\in S_3}\sigma(xyz),\\
 G_*={}&\Omega_1+
       u\Symm(xy)+v\Symm(yz)+(u+v)\Symm(zx)\\
 &+\frac1{18}\bigl(vx^2+(u+v)y^2+uz^2\bigr)\\
 &+\frac1{18}\bigl((-u^2-uv+3v^2)x
       +(3u^2+5uv+3v^2)y+(3u^2-uv-v^2)z\bigr).
\end{aligned}
\]
従って$B_*/(G_*+f)$ ($f\in k[u,v]_3$)は、
一つの周囲の中で相互に次数付き非同型な四次元代数族を与える。
さらに、$f$の非空Zariski開部分ではtailsが正則で大域次元3となる。
\end{proposition}
\begin{proof}
Fermat型の混合中心の普遍公式
\href{https://github.com/uedakazushi/nc-cubic/blob/main/papers/support/papers/sections/fermat_mixed_universal_center.tex}
{（nc-cubic論文A）}へ
$A=u,B=v,C=u+v,D=E=H=g=0,h=1$を代入すれば
上記の$G_*$とその中心性が得られる。
指定されたTHCRでは平行移動の位数が6なので
次数一・二の中心は零、次数三では$k\Omega_1$である。
相対PBWと中心元の三次還元から
$Z(B_*)_1=U$、$Z(B_*)_3=kG_*\oplus k[u,v]_3$となる。

任意の次数付き自己同型は中心一次空間を保つから、
$s=(u,v)^t$、$x=(x,y,z)^t$について
$s\mapsto Ns,\ x\mapsto Mx+\ell(s)$の形である。
命題\ref{prop:moduli-mixed-block-isomorphism}の勾配比較より
\[
 F_*(Mx+\ell(s),Ns)=\det(M)F_*(x,s)+r(s)
\]
となる。最高$x$-次数で
$\sum_i(Mx)_i^3=\det(M)\sum_ix_i^3$。
Fermat三次のHessianは非零定数倍の$xyz$なので、
$M$は座標三直線の合併を保ち、monomialである：
$(Mx)_i=\lambda_i x_{p(i)}$。
$x_{p(i)}^2s$項を比較すると$\ell=0$。
$\lambda_i^3=\det(M)$と
$\det(M)=\operatorname{sgn}(p)\prod_i\lambda_i$から
$p$は偶置換である。
同時スカラー倍で$\lambda_1=1$に規格化すると
$\lambda_i^3=1$、$\lambda_1\lambda_2\lambda_3=1$。
混合係数を
$(r_1,r_2,r_3)=(v,u+v,u)$とすれば
$r_i(Ns)=\lambda_i r_{p(i)}(s)$。
これらの一次式の唯一の関係の係数は
$c=(-1,1,-1)$である。よって
$\lambda_i=\mu c_{p(i)}/c_i$。
$p$が三cycleなら三つの比に$+1$と$-1$が混じり、
$\lambda_i^3=1$に矛盾する。
従って$p=1$、$\lambda_i=1$、$N=1$となり、
自己同型は同時スカラー倍だけ。
最後に次数二包絡・三次核の復元により、
$B_*/(G_*+f)\simeq_{\gr}B_*/(G_*+f')$なら
$[G_*+f]=[G_*+f']$。
$G_*$の係数を1と規格化したので$f=f'$である。
\end{proof}

\begin{corollary}[剛体周囲の中の正則な四係数非同型族]
\label{rem:moduli-rigid-generic-smoothness}
ある非零多項式$P(f)$が存在し、$P(f)\ne0$では
$\gldim\qgr B_*/(G_*+f)=3$である。
同時に可換な$V(F_*+f)\subset\PP^4$も滑らかとなる$f$は非空開部分をなし、
その中の有理係数点はZariski稠密である。
\end{corollary}
\begin{proof}
周囲はnoetherian Koszul CY5、相対associated gradedは$S_1\otimes k[u,v]$である。
第\ref{sec:pair-explicit-centers}節の中心式を代入すると$G_*$を得て、
全ての$f$で商は$k[u,v]$上自由である。
最高次中心元は非零正則、境界は位数6の滑らかな楕円THCRなので、
定理\ref{CSBGthm:global-open}の全仮定を満たす。
そこでは全係数chartの代数が$D_u[b,c,d]$となること、
中心持上げにより可換有限型中心上有限となることを確認してからBertiniを適用する。

可換な滑らかさについては$u=v=0$上に勾配の共通射影零点はない。
$u=1,t=v/u$上のactive勾配は
$x^2+y+(1+t)z$, $y^2+x+tz$, $z^2+(1+t)x+ty$であり、
$x,y,z$全次数を先に比較すると互いに素な先頭項$x^2,y^2,z^2$を持つ
monic Gr\"obner基底である。従って臨界スキームは$k[t]$上rank8の有限自由環を持ち、次元1である。
各点が全三様体の特異点となる条件は$(F_*+f)(P)=\partial_t(F_*+f)(P)=0$で、
$f(1,t)$の値と微分は四係数に対する独立な二条件である。
従って特異点incidenceの次元は高々3である。$v=1$も同様。
全incidenceは$\PP^4\times\A^4$内で閉で、その係数空間への像はproper性により閉である。
ゆえに滑らかな係数は非空開をなし、tails正則開部分とも交わる。
有理点の稠密性から最後の主張が従う。
これは可換三様体との残余圏同値を主張するものではない。
\end{proof}
